


\begin{comment}
Let $P(x)$ be a \emph{pencil} of geodesics based at $x$, i.e., the subset of $G^+(S)$ of all geodesics starting at $x \in \bdy\tilde{\Sigma}$.
By Proposition~\ref{prop:spikes}, the only way for a geodesic current $\mu$ to satisfy $\mu(P(z))>0$ for some $z\in \partial_{\infty} S$.
Further, note that if $\mu$ has no $0$-dimensional atoms, then it can't have any $1$-dimensional atoms, by the following result.

\begin{proposition}
  \label{prop:pencil}
  If $\mu(\widetilde{\gamma})=0$ for any $\widetilde{\gamma} \in G^+(S)$, then $\mu(P(z))=0$ for any $z \in \partial_{\infty} S$.
\end{proposition}
\begin{proof}
  Suppose by contradiction $\mu(P(z))>0$.  Then there must exist some $\widetilde{\gamma} \in P(z,[y_1,y_2]) \subset P(z)$ for some $[y_1,y_2] \subset \partial_{\infty} S$ disjoint from $z$, so that $p(\widetilde{\gamma})=\gamma$ is a closed geodesic. Let $g \in \pi_1(\Sigma)$ be the element so that its translation axis is $\widetilde{\gamma}$. The north-south dynamics of $\pi_1(\Sigma)$ on $\bdy\tilde{\Sigma}$ then forces $\cap_{n > 0} g^n P(a,[y_1,y_2]) = \tilde{\gamma}$. Note that $\mu(P(a,[y_1,y_2])) < \infty$ since $\mu$ is locally compact.  But then $0 \neq \mu(\cap_{n > 0} g^n P(z,[y_1,y_2]))=\mu(\tilde{\gamma})$, by $\pi_1(S)$-invariance of $\mu$, which is a contradiction, by Proposition~\ref{prop:spikes}.
\end{proof}
\end{comment}

%%% Local Variables:
%%% mode: latex
%%% TeX-master: "Intersections"
%%% End:
